\documentclass[12pt,a4paper,oneside]{article}
\usepackage[brazil]{babel}
\usepackage[utf8]{inputenc}
\usepackage{olymp}

\setcounter{problem}{4}

\begin{document}

\begin{problem}{Esparsidade}{esparsidade}{2}

Uma matriz é dita esparsa quando a ocorrência de elementos nulos (em geral o valor 0 para números) é frequente. Veja o exemplo a seguir:

$$A =
\begin{bmatrix}
0 & 23 & 0 & 0 & 0\\
0 & 2 & 0 & 0 & 0\\
7 & 0 & 0 & 0 & 1
\end{bmatrix}$$

A matriz $M$ possui alto índice de esparsidade, pois possui diversos elementos nulos. Numa matriz esparsa, pode-se ter muito tempo de processamento perdido com elementos nulos. Por exemplo, se quisermos somar os elementos da matriz $A$, precisamos apenas dos 4 elementos não nulos.
Uma forma de se obter maior desempenho é por meio da compactação de uma matriz esparsa. Pode-se compactar cada linha da matriz, armazenando-se apenas os elementos não nulos e a coluna em que aparecem. Uma representação compacta para a matriz $A$ seria:

$$A =
\begin{bmatrix}
(23,1)\\
(2,1)\\
(7,0)(1,4)
\end{bmatrix}$$

A taxa de compressão de uma matriz é dada pela razão entre o número de elementos na matriz após a compactação e o número total de elementos da matriz original. Para a matriz $A$ acima, a taxa de compressão é $tc=4/15=0.27$.

Você deve criar um programa que escreve a forma compactada de uma matriz esparsa bem como sua taxa de compressão.

%\Notes
%
%Observações, se tiver.

\InputFile


%A primeira linha contém um inteiro $N$, que indica o tamanho da matriz. Em seguida virão $N$ linhas com $N$ inteiros cada, indicando os elementos da matriz. Restrição: $1 \leq N \leq 1000$.

A primeira linha da entrada contém dois inteiros $M$ e $N$ indicando a dimensão da matriz. Em seguida são informados os valores da matriz $M \times N$, dados em $M$ linhas com $N$ inteiros cada. Restrições: $1 \leq M, N \leq 500$.


\OutputFile

O programa deverá escrever a matriz na forma compactada. Cada linha da matriz deve conter pares $(v,c)$ onde $v$ é um inteiro não nulo e $c$ é a coluna em que $v$ ocorre na linha. Cada par deve ser escrito entre parênteses e os elementos do par separados por
vírgula. Não deverá ser haver espaço entre pares. Linhas com todos os elementos nulos deverão conter um caractere `\#'.

Após escrever a matriz na forma compactada, o programa deve escrever o valor da taxa de compressão com duas casas decimais.

% \Notes
% Preste atenção aos tipos das variáveis, pois $F(90) \approx 2^{61}$.

\Example

\begin{example}
\exmp{
4 4
0 21 1 0
0 0 1 0
1 1 0 0
2 1 0 0
}{
(21,1)(1,2)
(1,2)
(1,0)(1,1)
(2,0)(1,1)
0.44
}%
\end{example}

\begin{example}
\exmp{
3 4
0 1 0 0
0 0 0 0
0 3 0 3
}{
(1,1)
\#
(3,1)(3,3)
0.25
}%
\end{example}

\begin{example}
\exmp{
3 3
0 0 0
0 0 0
0 0 0
}{
\#
\#
\#
0.00
}%
\end{example}

%Comentários sobre o exemplo, se necessário.
% Este exemplo corresponde ao da imagem mostrada acima.
\end{problem}

\end{document}
